\section{The unitary adjoint bound in the qualified quasiprimitive source class}
\label{sec:unitary-adjoint-quasiprimitive-corollary}

\paragraph{Accepted derived application and provenance.}
D-TD69Q, under question Q2fy, is an accepted derived application of
D-TD62 to the accepted standalone D-TD69 adjoint theorem. The parent
proposed the transfer; the complete proof is written by
\texttt{/root/sixfolds}. Whole author, independent and root reviews pass.
It remains separate from the immutable D-TD69 author freeze and retains
every written hypothesis of that theorem. No additional numerical
execution or separate new-obstruction credit is claimed. In particular it does not change the original field or
identify quasiprimitivity with primitivity.

Let \(N_*\) be the integer in D-TD62, (QPX.11), and put
\[
 M_0=\max\{80,N_*-2\},\qquad
 R_{\mathrm U}
   =\max_{\substack{q\text{ a prime power}\\2\le q<10^{12}}}
                 R(q,M_0),
 \tag{QUA.1}
\]
where \(R(q,M)\) is the explicitly computable integer in D-TD69,
(UAR.3). The finite set of prime powers is nonempty. Thus
\(R_{\mathrm U}\) is an explicitly computable integer; this assertion
does not claim that its relator-area algorithm has been executed.

\begin{corollary}[Qualified quasiprimitive unitary consequence]
Let an actual connected finite curve-cover family have a faithful
absolutely irreducible quasiprimitive coefficient representation \(V\),
with a model over its exact imaginary quadratic character field \(K\).
Assume its multiplicity variation has dimension six, signature
\((3,3)\), and actual period rank nine. Assume the naturally polarized
primitive sixfold cut out by a rational Rosati-fixed primitive
coefficient projector has nonsplit signed determinant type.

Assume the coefficient degree is minimal among all actual exact-\(K\)
sources with that signature, actual period rank and the same natural
signed determinant type. This comparison class includes nonprimitive
and nonquasiprimitive sources. Fix one complete recursion of
D-TD62, Sections~2--4.

Consider any nonabelian coarse layer of that recursion for which
every represented quasisimple component is an actual central quotient
of a transported named \(\operatorname{SU}_d(q)\) presentation.
Assume these components are quasisimple, as in the extraction theorem.
Then
\[
                         q<10^{12},\qquad d<R_{\mathrm U}.
 \tag{QUA.2}
\]
There is no character-level restriction and no assumed bound on its
component count. Its entire complementary factor is retained.
An arbitrary exceptional central cover is not included solely because
its simple quotient is unitary.
\end{corollary}

\begin{proof}
The accepted D-TD62 theorem applies to exactly the original family
and the stated comparison class. Its whole conclusion is used;
the following argument does not infer it from the primitive theorem.

\paragraph{The original complete floor.}
Sections~2--4 of D-TD62 give actual finite projective quotient maps
at every residual step and all separately stable coarse matrix
algebras. Their full tensor product is the original coefficient
space, by (QPX.7); the terminal one-dimensional multiplicity space
contributes only scalar matrices.

Take the subcollection of \emph{all} these original coarse layers
in (QPX.9). Its projective image is exactly the original projective
image on \(V\), not a new recursively decomposed complement.
The accepted alternating cutoff says that its largest alternating
natural degree \(m_*\) is strictly less than \(N_*\).
Integer degrees give \(m_*-1\le N_*-2\); the no-alternating
convention gives the same bound. Thus (QPX.9) supplies
\(\nu_V(T)/\dim V\ge1/M_0\) for every nonscalar original operator.
For a genuine scalar inertia \(zI\ne I\), the moving rank is
the full coefficient degree. Consequently, writing \(e=\dim V\),
every genuine original inertia satisfies
\[
                  \operatorname{rank}(T_i-I)\ge e/M_0.
 \tag{QUA.3}
\]
No lower bound is asserted for the number of nonidentity images
in any projected factor tuple.

\paragraph{The original sphere and twelve marks.}
The presence of a nontrivial quasisimple coarse layer makes \(e\ge2\).
For the original faithful nontrivial irreducible coefficient, let
\(h\) be the genus of its actual base, let \(r\) count precisely
the genuine original marks, and let \(D\) be the moving-rank sum.
The exact Euler and finite pointed-Hurwitz argument, written in
D-TD56 (NID.13), gives
\[
 6=2e(h-1)+D,\qquad r\le D,\qquad
 \operatorname{rank}(\mathrm{period})
       \le\dim\mathcal M_{h,r}.
 \tag{QUA.4}
\]
It uses the actual multiplicity variation and finite cover, not
primitivity of the coefficient representation.
If \(h=1\), then \(D=6\) and the pointed moduli dimension is at
most six, with the unpointed convention at most one.
If \(h=2\), then \(D=6-2e\ge0\) forces \(e\le3\) and the
dimension is at most \(3+D=9-2e\le5\).
For \(h\ge3\), \(D<0\). Actual period rank nine therefore forces
\(h=0\). The identity \(D=2e+6\) and the bound \(D\le re\)
give \(r>2\), so the rational pointed base is stable.
Its moduli dimension is \(r-3\); hence \(r\ge12\).
This proves exactly the original sphere, mark and budget hypotheses
of D-TD69. No family is replaced and no coefficient-field extension
or polarization comparison is made in this step.

\paragraph{The actual selected layer.}
The chosen nonabelian extraction has minimal normal projective core
\(S^t\). D-TD62 Section~3 constructs its inverse-image derived
groups and proves that they are commuting quasisimple components
with central intersections. They act faithfully and irreducibly on
the extracted whole \(W\), and its factorization into faithful
irreducible component modules is exactly (QPX.6).
Minimal normality makes the component action transitive.

The coarse layer \(W\) is separately stable under every original
operator. Equations (QPX.1)--(QPX.2), pulled back through all earlier
residual maps, give an actual factorization
\(V=W\otimes U\) and an actual finite projective factor image.
It contains the indicated represented core and normalizes its
projective image. The transported named presentation is the same
on every component. These are precisely the standalone D-TD69
layer hypotheses. No primitivity of this coarse image or of an
arbitrary whole complement is required. Nor is any later core
asserted to lift as a subgroup of the original coefficient group.

\paragraph{The high-field restriction.}
The accepted D-TD03 applies directly to this actual \(W\), with
the entire original \(U\). Indeed the finite special-linear lift
of its projective image contains the represented core by perfectness;
taking derived groups upgrades projective normalization to ordinary
normalization. This is the full-normalizer interface proved in
D-TD03 itself. It does not require the original-subgroup premise
of D-TD04 for a later extracted core.

If \(q\ge10^{12}\), all selected components have large defining
field in their transported named presentation, and their simple
quotients are nonalternating. The last assertion follows from the
full finite-simple-group repetition list already used in D-TD68;
all unitary alternating coincidences lie at small field and dimension.
D-TD03 then excludes the actual source. Thus \(q<10^{12}\).
No mixed small/large-field core is assigned its uniform trace bound.

\paragraph{Apply the standalone adjoint theorem.}
If \(d\ge120\), all D-TD69 hypotheses have now been supplied:
(QUA.3), the original sphere and at least twelve marks, its exact
moving sum, and the actual transitive named unitary coarse layer
with the entire partner. That theorem gives
\(d<R(q,M_0)\le R_{\mathrm U}\).
If \(d<120\), the same conclusion is automatic because the
definition (UAR.3) has \(R(q,M_0)\ge4\cdot10^{33}\).
This proves (QUA.2).
\end{proof}

\paragraph{Exact boundary of the transfer.}
This corollary uses the full accepted quasiprimitive D-TD62 theorem
with its original exact field, actual rank-nine family, natural
signed determinant and all-source minimality comparison.
It does not assert that a general minimal source is quasiprimitive.
The proper-field canonical-induction cases left by D-TD56/60/67
remain separate. It adds no field descent, polarization change,
actual-source construction, arbitrary original-group order bound
or algebraic cycle. The standalone D-TD69 mathematical proof and its
primitive-source application retain their separately reviewed scopes.
